\section{Sets and classes}
\label{sec:set-prerequisites-01}

\begin{definition}[Subset]
For sets \(A\) and \(B\), write \(A\subseteq B\) iff
\[
  \forall x\,(x\in A\to x\in B).
\]
The sets are equal iff each is a subset of the other.
\end{definition}

\begin{definition}[Power set]
For a set \(A\), its \emph{power set} is
\[
  \Pow(A)\defeq\set{X:X\subseteq A}.
\]
\end{definition}

\begin{definition}[Class]
A \emph{class} is a metalinguistically specified collection of sets.  A class
need not be a set.  A \emph{proper class} is a class that is not a set.
\end{definition}

An unrestricted collection described by a condition may be a proper class.
Quantification over such a collection is metalinguistic quantification in the
ambient foundation, not object-language quantification over a set.

\begin{warning}
Set-builder notation does not itself prove that a set exists.  The expression
\(\set{x:P(x)}\) may denote a proper class.  When existence matters, identify a
previously available bounding set or cite the relevant set-existence axiom.
\end{warning}

